\chapter{Conclusion}
\label{chap:conclusion}

We asked whether sparse autoencoders can identify human-interpretable features
in a language model's residual stream that explain and causally influence
whether the model adopts or resists feedback. For a Gemma~3 4B Solver
answering multi-hop questions with feedback from a Qwen3 4B Critic, the answer
is: partly.

The Solver's decision is readable at layer~25 once it has seen the feedback,
and SAE features carry much of that signal, although less than the raw
activations. Editing the selected features of the pretrained SAE family
changed the decision more often than equal-size random edits, in four of nine
prespecified tests, by 1.6 to 2.5 percentage points. The effect is small,
mostly pushes the Solver from adopting to retaining, and in a post hoc split
appeared on the natural MuSiQue questions rather than the synthetic ones. The features are
not, however, interpretable as feedback uptake itself. They mostly encode the
kind of answer the Critic gave, and the features that were easiest to describe
were the weakest causal handles. SAEs in this setting reveal variables the
decision depends on, not the decision.

The study's main contribution is a pipeline and an evaluation framework that
keep these kinds of evidence apart: reconstruction, prediction, meaning and
causal effect, each tested against matched controls. The earlier designs,
whose effects disappeared once reconstruction, random and equal-size controls
were added, show why each control is needed. A feature that predicts a
behavior, or one that can steer it, is not yet an explanation of it.
